\xchapter{The Teaching of the Twelve Apostles}

\xsection{Chapter I.—The Two Ways; The First Commandment}

1. \textsc{There} are two ways,\footnote{\par  This phrase connects the book with the Duæ Viæ; see Introductory Notice. Barnabas has “light” and “darkness” for “life” and “death.”} one of life and one of death;\footnote{\par  Deut. xxx. 15, 19; Jer. xxi. 8; Matt. vii. 13, 14} but a great difference between the two ways. 2. The way of life, then, is this: First, thou shalt love God\footnote{\par  Comp. Deut. vi. 5, which is fully cited in Apostolic Constitutions, vii. 2, though the verb here is more exactly cited from LXX.} who made thee; second, thy neighbour as thyself;\footnote{\par  Lev. xix. 18; Matt. xxii. 37, 39. Comp. Mark xii. 30, 31} and all things whatsoever thou wouldst should not occur to thee, thou also to another do not do.\footnote{\par  Comp. Tobit iv. 15; and Matt. vii. 12; Luke vi. 31} 3. And of these sayings\footnote{\par  These Old-Testament commands are thus taught by the Lord.} the teaching is this: Bless them that curse you, and pray for your enemies, and fast for them that persecute you.\footnote{\par  Matt. v. 44. But the last clause is added, and is of unknown origin; not found in Apostolic Constitutions} For what thank is there, if ye love them that love you? Do not also the Gentiles do the same?\footnote{\par  Matt. v. 46, 47; Luke vi. 32. The two passages are combined.} But do ye love them that hate you; and ye shall not have an enemy.\footnote{\par  So Apostolic Constitutions. Comp. 1 Pet. iii. 13} 4. Abstain thou from fleshly and worldly lusts.\footnote{\par  1 Pet. ii. 11. The Codex has \textgreek{σωματικῶν}, “bodily;” but editors correct to \textgreek{κοσμικῶν}} If one give thee a blow upon thy right cheek, turn to him the other also;\footnote{\par  Matt. v. 39; Luke vi. 29.} and thou shalt be perfect. If one impress thee for one mile, go with him two.\footnote{\par  Matt. v. 41} If one take away thy cloak, give him also thy coat.\footnote{\par  Matt. v. 40; Luke vi. 29} If one take from thee thine own, ask it not back,\footnote{\par  Luke vi. 30. The last clause is a peculiar addition: “art not able,” since thou art a Christian; otherwise it is a commonplace observation.} for indeed thou art not able. 5. Give to every one that asketh thee, and ask it not back;\footnote{\par  Luke vi. 30. The rest of the sentence is explained by the parallel passage in Apostolic Constitutions, which cites Matt. v. 45.} for the Father willeth that to all should be given of our own blessings (free gifts).\footnote{\par  Bryennios finds a parallel (or citation) in Hermas, Commandment Second, p. 20, vol. i. Ante-Nicene Fathers. The remainder of this chapter has no parallel in Apostolic Constitutions.} Happy is he that giveth according to the commandment; for he is guiltless. Woe to him that receiveth; for if one having need receiveth, he is guiltless; but he that receiveth not having need, shall pay the penalty, why he received and for what, and, coming into straits (confinement),\footnote{\par  Gr. \textgreek{ἐν συνοχῇ}. Probably = imprisonment; see next clause.} he shall be examined concerning the things which he hath done, and he shall not escape thence until he pay back the last farthing.\footnote{\par  Matt. v. 26.} 6. But also now concerning this, it hath been said, Let thine alms sweat\footnote{\par  Codex: \textgreek{ιδροτάτω}, which in this connection is unintelligible. Bryennios corrects into \textgreek{ιδροσάτω}, rendered as above. There are various other conjectural emendations. The verse probably forbids indiscriminate charity, pointing to an early abuse of Christian liberality.} in thy hands, until thou know to whom thou shouldst give.

\xsection{Chapter II.—The Second Commandment: Gross Sin Forbidden}

1. And the second commandment of the Teaching; 2. Thou shalt not commit murder, thou shalt not commit adultery,\footnote{\par  Ex. xx. 13, 14.} thou shalt not commit pæderasty,\footnote{\par  Or, “corrupt boys,” as in the version of Apostolic Constitutions.} thou shalt not commit fornication, thou shalt not steal,\footnote{\par  Ex. xx. 15.} thou shalt not practice magic, thou shalt not practice witchcraft, thou shalt not murder a child by abortion nor kill that which is begotten.\footnote{\par  Comp. Ex. xxi. 22, 23. The Codex reads \textgreek{γεννηθέντα}, which Schaff renders “the new-born child.” Bryennios substitutes \textgreek{γεννηθέν}, which is accepted by most editors, and rendered as above.} Thou shalt not covet the things of thy neighbour,\footnote{\par  Ex. xx. 17.} 3. thou shalt not forswear thyself,\footnote{\par  Matt. v. 34.} thou shalt not bear false witness,\footnote{\par  Ex. xx. 16.} thou shalt not speak evil, thou shalt bear no grudge.\footnote{\par  Rendered “nor shalt thou be mindful of injuries” in version of Apostolic Constitutions.} 4. Thou shalt not be double-minded nor double-tongued; for to be double-tongued is a  snare of death.\footnote{\par  So Barnabas, xix.} 5. Thy speech shall not be false, nor empty, but fulfilled by deed.\footnote{\par  Verse 5, except the first clause, occurs only here.} 6. Thou shalt not be covetous, nor rapacious, nor a hypocrite, nor evil disposed, nor haughty. Thou shalt not take evil counsel against thy neighbour.\footnote{\par  Latter half of verse 6 in Barnabas, xix.} 7. Thou shalt not hate any man; but some thou shalt reprove,\footnote{\par  Lev. xix. 17; Apostolic Constitutions.} and concerning some thou shalt pray, and some thou shalt love more than thy own life.\footnote{\par  Or, “soul.” The last part of the clause is found in Barnabas; but “and concerning some…pray, and some” has no parallel. An interesting verse in its literary history.}

\xsection{Chapter III.—Other Sins Forbidden}

1. My child,\footnote{\par  The address “my child” does not occur in the parallel passages.} flee from every evil thing, and from every likeness of it. 2. Be not prone to anger, for anger leadeth the way to murder; neither jealous, nor quarrelsome, nor of hot temper; for out of all these murders are engendered. 3. My child, be not a lustful one; for lust leadeth the way to fornication; neither a filthy talker, nor of lofty eye; for out of all these adulteries are engendered. 4. My child, be not an observer of omens, since it leadeth the way to idolatry; neither an enchanter, nor an astrologer, nor a purifier, nor be willing to took at these things; for out of all these idolatry is engendered. 5. My child, be not a liar, since a lie leadeth the way to theft; neither money-loving, nor vainglorious, for out of all these thefts are engendered. 6. My child, be not a murmurer, since it leadeth the way to blasphemy; neither self-willed nor evil-minded, for out of all these blasphemies are engendered. 7. But be thou meek, since the meek shall inherit the earth.\footnote{\par  Matt. v. 5.} 8. Be long-suffering and pitiful and guileless and gentle and good and always trembling at the words which thou hast heard.\footnote{\par  Isa. lxvi. 2, 5; Apostolic Constitutions, vii. 8.} 9. Thou shalt not exalt thyself,\footnote{\par  Comp. Luke xviii. 14.} nor give over-confidence to thy soul. Thy soul shall not be joined with lofty ones, but with just and lowly ones shall it have its intercourse. 10. The workings that befall thee receive as good, knowing that apart from God nothing cometh to pass.\footnote{\par  Ecclus. ii. 4. So Bryennios. Comp. last part of Apostolic Constitutions vii. 8.}

\xsection{Chapter IV.—Various Precepts}

1. My child, him that speaketh to thee the word of God remember night and day; and thou shalt honour him as the Lord;\footnote{\par  Comp. Heb. xiii. 7. In Apostolic Constitutions there is a transposition of words.} for in the place whence lordly rule is uttered,\footnote{\par  Schaff: “The Lordship is spoken of.” Apostolic Constitutions, “where the doctrine concerning God is,” etc.} there is the Lord. 2. And thou shalt seek out day by day the faces of the saints, in order that thou mayest rest upon\footnote{\par  Or, “acquiesce in” (Apostolic Constitutions).} their words. 3. Thou shalt not long for\footnote{\par  Some read \textgreek{ποιήσεις}, “make,” as in Apostolic Constitutions and Barnabas, instead of \textgreek{ποθήσεις}, Codex.} division, but shalt bring those who contend to peace. Thou shalt judge righteously, thou shalt not respect persons in reproving for transgressions. 4. Thou shalt not be undecided whether it shall be or no.\footnote{\par  Comp. Ecclus. i. 28. The verse occurs in Barnabas; and in Apostolic Constitutions “in thy prayer” is inserted, which is probably the sense here.} 5. Be not a stretcher forth of the hands to receive and a drawer of them back to give.\footnote{\par  Ecclus. iv. 31. The Greek word \textgreek{συσπῶν} occurs here and in Barnabas, but not in Apostolic Constitutions.} 6. If thou hast aught, through thy hands thou shalt give ransom for thy sins.\footnote{\par  Apostolic Constitutions adds, in explanation, Prov. xvi. 6.} 7. Thou shalt not hesitate to give, nor murmur when thou givest; for thou shalt know who is the good repayer of the hire. 8. Thou shalt not turn away from him that is in want, but thou shalt share all things with thy brother, and shalt not say that they are thine own; for if ye are partakers in that which is immortal, how much more in things which are mortal?\footnote{\par  Comp. Acts iv. 32; Rom. xv. 27. The latter half of the verse is in Barnabas (not in Apostolic Constitutions), but with the substitution of “incorruptible” and “corruptible.”} 9. Thou shalt not remove thy hand from thy son or from thy daughter, but from their youth shalt teach them the fear of God.\footnote{\par  Comp. Eph. vi. 4.} 10. Thou shalt not enjoin aught in thy bitterness upon thy bondman or maidservant, who hope in the same God, lest ever they shall fear not God who is over both;\footnote{\par  Comp. Eph. vi. 9; Col. iv. 1.} for he cometh not to call according to the outward appearance, but unto them whom the Spirit hath prepared. 11. And ye bondmen shall be subject to your\footnote{\par  Codex reads “our;” editors correct to “your. ”} masters as to a type of God, in modesty and fear.\footnote{\par  Comp. Eph. vi. 5; Col. iii. 22.} 12. Thou shalt hate all hypocrisy and everything which is not pleasing to the Lord. 13. Do thou in no wise forsake the commandments of the Lord; but thou shalt keep what thou hast received, neither adding thereto nor taking away therefrom.\footnote{\par  Deut. xii. 32.} 14. In the church\footnote{\par  “In the congregation;” i.e., assembly of believers. This phrase is omitted in both Barnabas and Apostolic Constitutions. Comp. Jas. v. 16.} thou shalt acknowledge thy transgressions, and thou shalt not come near for thy prayer\footnote{\par  Or, “to thy place of prayer” (Schaff).} with an evil conscience.\footnote{\par  So Barnabas; but Apostolic Constitutions, “in the day of thy bitterness.”} This is the way of life.\footnote{\par  So Apostolic Constitutions; but Barnabas, “the way of light.” See note on chap. i. 1.}

\xsection{Chapter V.—The Way of Death}

1. And the way of death\footnote{\par  Barnabas has “darkness,” but afterwards “way of eternal death.”} is this: First of all it is evil and full of curse:\footnote{\par  Not in Apostolic Constitutions, and no exact parallel in Barnabas.} murders,\footnote{\par  Of the twenty-two sins named in this verse, Barnabas gives fourteen, in differing order, and in the singular; Apostolic Constitutions gives all but one (\textgreek{υψος}, “loftiness” “haughtiness”), in the same order, and with the same change from plural to singular.} adulteries, lusts, fornications, thefts, idolatries, magic arts, witchcrafts, rapines, false witnessings, hypocrisies, double-heartedness, deceit, haughtiness, depravity, self-will, greediness, filthy talking, jealousy, over-confidence, loftiness, boastfulness; 2. persecutors of the good,\footnote{\par  This verse appears almost word for word in Barnabas, with two additional clauses.} hating truth, loving a lie, not knowing a reward for righteousness, not cleaving\footnote{\par  The Apostolic Constitutions give a parallel from this point; verbally exact from the phrase, “not for that which is good.”} to good nor to righteous judgment, watching not for that which is good, but for that which is evil; from whom meekness and endurance are far, loving vanities, pursuing requital, not pitying a poor man, not labouring for the afflicted, not knowing Him that made them, murderers of children, destroyers of the handiwork of God, turning away from him that is in want, afflicting him that is distressed, advocates of the rich, lawless judges of the poor, utter sinners.\footnote{\par  The word \textgreek{πανθαμαρτητοι} occurs only here, and in the parallel passage in Barnabas (rendered in this edition “who are in every respect transgressors,” vol. i. p. 149), and in Apostolic Constitutions (rendered “full of sin”). A similar term occurs in the recently recovered portion of 2 Clement, xviii., where Bishop Lightfoot renders, as above, “an utter sinner.”} Be delivered, children, from all these.\footnote{\par  Found verbatim in Apostolic Constitutions, not in Barnabas: with the latter there is no further parallel, except a few phrases in chap. xvi. 2, 3 (which see).}

\xsection{Chapter VI.—Against False Teachers, and Food Offered to Idols}

1. See that no one cause thee to err\footnote{\par  Comp. Matt. xxiv . 4 (Greek); Revised Version, “lead you astray:” Apostolic Constitutions, vii. 19.} from this way of the Teaching, since apart from God it teacheth thee. 2. For if thou art able to bear all the yoke\footnote{\par  Or, “the whole yoke.” Those who accept the Jewish-Christian authorship refer this to the ceremonial law. It seems quite as likely to mean ascetic regulations. Of these there are many traces, even in the New-Testament churches.} of the Lord, thou wilt be perfect; but if thou art not able, what thou art able that do. 3. And concerning food,\footnote{\par  Apostolic Constitutions, vii. 20, begins with a similar phrase, but is explicitly against asceticism in this respect. The precepts here do not indicate any such spirit as that opposed by Paul.} bear what thou art able; but against that which is sacrificed to idols\footnote{\par  Comp. Acts xv. 20, 29; 1 Cor. viii. 4, etc., x. 18, etc. (Rom. xiv. 20 refers to ascetic abstinence.) This prohibition had a necessary permanence; comp. Apostolic Constitutions, vii. 21.} be exceedingly on thy guard; for it is the service of dead gods.\footnote{\par  Comp. the same phrase in 2 Clement, iii. This chapter closes the first part of the Teaching, that supposed to be intended for catechumens. The absence of doctrinal statement does not necessarily prove the existence of a circle of Gentile Christians where the Pauline theology was unknown. If such a circle existed, emphasizing the ethical side of Christianity to the exclusion of its doctrinal basis, it disappeared very soon. From the nature of the case, that kind of Christianity is intellectually weak and necessarily short-lived.}

\xsection{Chapter VII.—Concerning Baptism}

1. And concerning baptism,\footnote{\par  Verse 1 is found, well-nigh entire, in Apostolic Constitutions vii. 22, but besides this only a few words of verses 2 and 4. The chapter has naturally called out much discussion as to the mode of baptism.} thus baptize ye:\footnote{\par  [Elucidation I.]} Having first said all these things, baptize into the name of the Father, and of the Son, and of the Holy Spirit,\footnote{\par  Matt. xxviii. 19.} in living water.\footnote{\par  Probably running water.} 2. But if thou have not living water, baptize into other water; and if thou canst not in cold, in warm. 3. But if thou have not either, pour out water thrice\footnote{\par  The previous verses point to immersion; this permits pouring in certain cases, which indicates that this mode was not unknown. The trine application of the water, and its being poured on the head, are both significant.} upon the head into the name of Father and Son and Holy Spirit. 4. But before the baptism let the baptizer fast, and the baptized, and whatever others can; but thou shalt order the baptized to fast one or two days before.\footnote{\par  The fasting of the baptized is enjoined in Apostolic Constitutions, but that of the baptizer (and others) is peculiar to this document.}

\xsection{Chapter VIII.—Concerning Fasting and Prayer (the Lord’s Prayer)}

1. But let not your fasts be with the hypocrites;\footnote{\par  Comp. Matt. vi. 16.} for they fast on the second and fifth day of the week; but do ye fast on the fourth day and the Preparation (Friday).\footnote{\par  The reasons for fasting on Wednesday and Friday are given in Apostolic Constitutions (the days of betrayal and of burial). Monday and Thursday were the Jewish fast-days. The word “Preparation” (day before the Jewish sabbath) occurs in Matt. xxvii. 62, etc., and for some time retained a place in Christian literature.} 2. Neither pray as the hypocrites; but as the Lord commanded in His Gospel,\footnote{\par  Matt. vi. 5, 9–13. This form of the Lord’s Prayer is evidently cited from Matthew, not from Luke. The textual variations are slight. The citation is of importance as proving that the writer used this Gospel, and that the liturgical use of the Lord’s Prayer was common.} thus pray: Our Father who art in heaven, hallowed be Thy name. Thy kingdom come. Thy will be done, as in heaven, so on earth. Give us to-day our daily (needful) bread,\footnote{\par  On this phrase, comp. Revised Version, Matt. vi. 11; Luke xi. 3 (text, margin, and American appendix).} and forgive us our debt as we also forgive our debtors. And bring us not into temptation, but deliver us from the evil one (or, evil); for Thine is the power and the glory for ever.\footnote{\par  The variation in the form of the doxology confirms the judgment of textual criticism, which omits it in Matt. vi. 13. All early liturgical literature tends in the same direction; comp. Apostolic Constitutions, vii. 24.} 3. Thrice in the day thus pray.\footnote{\par  This is in accordance with Jewish usage. Dan. vi. 10; Ps. lv. 17. Comp. Acts iii. 1, x. 9.}

\xsection{Chapter IX.—The Thanksgiving (Eucharist)}

1. Now concerning the Thanksgiving (Eucharist), thus give thanks. 2. First, concerning the cup:\footnote{\par  This is a variation from the order of the New Testament and of all liturgies: probably this led to its omission in Apostolic Constitutions. The word “for” may be substituted for “concerning” here and in verse 3. [Possibly a response for recipients.]} We thank thee, our Father, for the holy vine of David Thy servant,\footnote{\par  Peculiar to this passage, but derived from a common scriptural figure and from the paschal formula. Comp. especially John xv. 1; Matt. xxvi. 29; Mark xiv. 25.} which Thou madest known to us through Jesus Thy Servant; to Thee be the glory for ever. 3. And concerning the broken bread:\footnote{\par  The word \textgreek{κλάσμα} is found in the accounts of the feeding of the multitude (Matt. xiv. 20, xv. 37, and parallels); it was naturally applied to the broken bread of the Eucharist.} We thank Thee, our Father, for the life and knowledge which Thou madest known to us through Jesus Thy Servant; to Thee be the glory for ever. 4. Even as this broken bread was scattered over the hills,\footnote{\par  This reference to “hills,” or “mountains,” is used as an argument against the Egyptian origin of the Teaching.} and was gathered together and became one, so let Thy Church be gathered together from the ends of the earth into Thy kingdom;\footnote{\par  This part of the verse is found in Apostolic Constitutions. Schaff properly calls attention to the distinction here made between “Thy Church” and “Thy kingdom.”} for Thine is the glory and the power through Jesus Christ for ever. 5. But let no one eat or drink of your Thanksgiving (Eucharist), but they who have been baptized into the name of the Lord; for concerning this also the Lord hath said, Give not that which is holy to the dogs.\footnote{\par  Matt. vii. 6.}

\xsection{Chapter X.—Prayer After Communion}

1. But after ye are filled,\footnote{\par  “After the participation” (Apostolic Constitutions) points to a distinct Eucharistic service. Here the Lord’s Supper is evidently connected with the Agape [a noteworthy suggestion]; comp. 1 Cor. xi. 20–22, 33. This is an evidence of early date; comp. Justin Martyr, Apol. i. chaps. 64–66, where the Lord’s Supper is shown to be distinct (Ante-Nicene Fathers, i. pp. 185, 186).} thus give thanks: 2. We thank Thee, holy Father, for Thy holy name which Thou didst cause to tabernacle in our hearts, and for the knowledge and faith and immortality, which Thou madest known to us through Jesus Thy Servant; to Thee be the glory for ever. 3. Thou, Master almighty, didst create all things for Thy name’s sake; Thou gavest food and drink to men for enjoyment, that they might give thanks to Thee; but to us Thou didst freely give spiritual food and drink and life eternal through Thy Servant.\footnote{\par  This last clause has no parallel in Apostolic Constitutions, and points to an earlier and more spiritual conception of the Eucharist. Verse 4 also is peculiar to this passage.} 4. Before all things we thank Thee that Thou art mighty; to Thee be the glory for ever. 5. Remember, Lord, Thy Church, to deliver it from all evil and to make it perfect in Thy love, and gather it from the four winds, sanctified for Thy kingdom which Thou hast prepared for it;\footnote{\par  The above rendering follows Bryennios; that of Harnack (formerly accepted by Hall and Napier) is: “Gather it, sanctified, from the four winds, into Thy kingdom,” etc. The phrase “from the four winds” recalls Matt. xxiv. 31.} for Thine is the power and the glory for ever. 6. Let grace come, and let this world pass away.\footnote{\par  This is peculiar; but comp. 1 Cor. vii. 31 for the last clause.} Hosanna to the God (Son)\footnote{\par  The Codex reads \textgreek{τῷ ὺεῷ}, which Bryennios alters to \textgreek{τῷ ὺιῷ}. The former is the more difficult reading, and is defended by Harnack.} of David! If any one is holy, let him come; if any one is not so, let him repent.\footnote{\par  This exhortation indicates a mixed assembly; comp. Apostolic Constitutions. [If so, it belongs to the Agape.]} Maran atha.\footnote{\par  Cor. xvi. 22, Revised Version, margin: “That is, our Lord cometh.” Comp. Rev. xxii. 20.} Amen. 7. But permit the prophets to make Thanksgiving as much as they desire.\footnote{\par  A limitation as compared with 1 Cor. Xiv. 29, 31, and yet indicating a combination of extemporaneous devotion with the liturgical form. The verse prepares the way for the next chapter.}

\xsection{Chapter XI.—Concerning Teachers, Apostles, and Prophets}

1. Whosoever, therefore, cometh and teacheth you all these things that have been said before, receive him.\footnote{\par  This refers to all teachers, more fully described afterwards.} 2. But if the teacher himself turn\footnote{\par  Lit. “being turned:” i.e. turned from the truth, perverted.} and teach another doctrine to the destruction of this, hear him not; but if he teach so as to increase righteousness and the knowledge of the Lord, receive him as the Lord. 3. But concerning the apostles and prophets, according to the decree of the Gospel, thus do. 4. Let every apostle that cometh to you be received as the Lord.\footnote{\par  Matt. x. 40. The mention of apostles here has caused much discussion, but there are many indications that travelling evangelists were thus termed for some time after the apostolic age. Bishop Lightfoot has shown, that, even in the New Testament, a looser use of the term applied it to others than the Twelve. Comp. Rom. xvi. 7; 1 Cor. xv. 5, 7 (?); Gal. i. 19; 1 Thess. ii. 6: also, as applied to Barnabas, Acts xiv. 4, 14.} 5. But he shall not remain except one day; but if there be need, also the next; but if he remain three days, he is a false prophet. 6. And when the apostle goeth away, let him take nothing but bread until he lodgeth;\footnote{\par  Reach a place where he can lodge.} but if he ask money, he is a false prophet. 7. And every prophet that speaketh in the Spirit\footnote{\par  Under the influence of the charismatic gift spoken of in 1 Cor. xii. 3, xiv. 2. Another indication of an early date.} ye shall neither try nor judge; for every sin shall be forgiven, but this sin shall not be forgiven.\footnote{\par  Probably a reference to the sin against the Holy Spirit. Matt. xii. 31, 32; Mark iii. 29, 30.} 8. But not every one that speaketh in the Spirit is a prophet; but only if he hold the ways of the Lord. Therefore from their ways shall the false prophet and the prophet be known. 9. And every prophet who ordereth a meal\footnote{\par  Probably a love-feast, commanded by the prophet in his peculiar utterance.} in the Spirit eateth not from it, except indeed he be a false prophet; 10. and every prophet who teacheth the truth, if he do not what he teacheth, is a false prophet. 11. And every prophet, proved true,\footnote{\par  \textgreek{ἀληθινός}, “genuine.”} working unto the mystery of the Church in the world,\footnote{\par  \textgreek{ποιῶν εἰς μυστήριον κοσμικὸν ἐκκλησίας,} “working unto a worldly mystery of (the) Church,” or “making assemblies for a worldly mystery.” Either rendering is grammatical: neither is very intelligible. The paraphrase in the above version presents one leading view of this difficult passage: the mystery is the Church, and a worldly one, because the Church is in the world. The other leading view joins \textgreek{ἐκκλησίας} (as accusative) with \textgreek{ποιῶν}, “making assemblies for a worldly mystery.” So Bryennios, who regards the worldly mystery as a symbolical act of the prophet. Others suggest, as the mystery for which the assemblies are called, revelation of future events, celibacy, the Eucharist, the ceremonial law. It seems, at all events, to point to incipient fanaticism on the part of the prophets of those days. [Elucidation III.] This was likely to take the form either of asceticism or of extravagant predictions and mystical fancies about the Church in the world. Did we know the place and the time more accurately, we might decide which was meant. This caution was evidently needed: Let God judge such extravagances.} yet not teaching others to  do what he himself doeth, shall not be judged among you, for with God he hath his judgment; for so did also the ancient prophets. But whoever saith in the Spirit, Give me money, or something else, ye shall not listen to him; but if he saith to you to give for others’ sake who are in need, let no one judge him.

\xsection{Chapter XII.—Reception of Christians}

1. But let every one that cometh in the name of the Lord be received,\footnote{\par  All professed Christians are meant.} and afterward ye shall prove and know him; for ye shall have understanding right and left. 2. If he who cometh is a wayfarer, assist him as far as ye are able; but he shall not remain with you, except for two or three days, if need be. 3. But if he willeth to abide with you, being an artisan, let him work and eat;\footnote{\par  Comp. 2 Thess. iii. 10.} but if he hath no trade, 4. according to your understanding see to it that, as a Christian,\footnote{\par  The term occurs only here in the Teaching.} he shall not live with you idle. 5. But if he willeth not to do, he is a Christ-monger.\footnote{\par  “Christ-trafficker.” The abuse of Christian fellowship and hospitality naturally followed the remarkable extension of Christianity. This expressive term was coined to designate the class of idlers who would make gain out of their professed Christianity. It occurs in the longer form of the Ignatian Epistles (Trallians, vi.) and in literature of the fourth century.} Watch that ye keep aloof from such.

\xsection{Chapter XIII.—Support of Prophets}

1. But every true prophet that willeth to abide among you\footnote{\par  “Who will settle among you” (Hitchcock and Brown). The itinerant prophets might become stationary, we infer. Chaps. xi.-xv. point to a movement from an itinerant and extraordinary ministry to a more settled one.} is worthy of his support.\footnote{\par  Lit., “nourishment,” “food.”} 2. So also a true teacher is himself worthy, as the workman, of his support.\footnote{\par  Matt. x. 10; comp. Luke x. 7.} 3. Every first-fruit, therefore, of the products of wine-press and threshing-floor, of oxen and of sheep, thou shalt take and give to the prophets, for they are your high priests.\footnote{\par  This phrase, indicating a sacerdotal view of the ministry, seems to point to a later date than that claimed for the Teaching. Some regard it as an interpolation: others take it in a figurative sense. In Apostolic Constitutions the sacerdotal view is more marked. [1 Pet. ii. 9. If the plebs = “priests,” prophets = “high priests.”] Here the term is restricted to the prophets: compare Schaff in loco.} 4. But if ye have not a prophet, give it to the poor. 5. If thou makest a batch of dough, take the first-fruit and give according to the commandment. 6. So also when thou openest a jar of wine or of oil, take the first-fruit and give it to the prophets; 7. and of money (silver) and clothing and every possession, take the first-fruit, as it may seem good to thee, and give according to the commandment.

\xsection{Chapter XIV.—Christian Assembly on the Lord’s Day}

1. But every Lord’s day\footnote{\par  Comp. Rev. i. 10. Here the full form is \textgreek{κατὰ κυριακὴν δὲ Κυρίον}. If the early date is allowed, this verse confirms the view that from the first the Lord’s Day was observed, and that, too, by a eucharistic celebration.} do ye gather yourselves together, and break bread, and give thanksgiving after having confessed your transgressions,\footnote{\par  Comp chap. iv. 14. No parallel in Apostolic Constitutions.} that your sacrifice may be pure.\footnote{\par  On this spiritual sense of “sacrifice,” comp. Rom. xii. 1; Phil. ii. 17; Heb. xiii. 15; 1 Pet. ii. 5.} 2. But let no one that is at variance\footnote{\par  “That hath the (or, any) dispute” (\textgreek{ἀμφιβολίαν}); comp. Matt. v. 23, 24.} with his fellow come together with you, until they be reconciled, that your sacrifice may not be profaned. 3. For this is that which was spoken by the Lord: In every place and time offer to me a pure sacrifice;\footnote{\par  [See Mal. i. 11. See Irenæus, cap. xvii. 5, vol. i. p. 484.]} for I am a great King, saith the Lord, and my name is wonderful among the nations.\footnote{\par  Mal. i. 11, 14. Quoted in Apostolic Constitutions and by several Ante-Nicene Fathers, with the same reference to the Eucharist.}

\xsection{Chapter XV.—Bishops and Deacons; Christian Reproof}

1. Appoint, therefore, for yourselves, bishops and deacons worthy of the Lord, men meek, and not lovers of money,\footnote{\par  Comp. 1 Tim. iii. 4.} and truthful and proved; for they also render to you the service\footnote{\par  Or, “ministry.” This clause and the following verse indicate that the extraordinary ministers were as yet more highly regarded.} of prophets and teachers. 2. Despise them not therefore, for they are your honoured ones, together with the prophets and teachers. 3. And reprove one another, not in anger, but in peace, as ye have it in the Gospel;\footnote{\par  Comp. Matt. xviii. 15–17.} but to every one that acts amiss\footnote{\par  The word \textgreek{ἀστοχέω}, occurring here, means “to miss the mark;” in New Testament, “to err” or, “swerve.” See 1 Tim. i. 6, vi. 21; 2 Tim. ii. 18.} against another, let no one speak, nor let him hear aught from you until he repent. 4. But your prayers and alms and all your deeds so do, as ye have it in the Gospel of our Lord.\footnote{\par  The reference here is probably to the Sermon on the Mount: Matt. v.-vii., especially to chap. vi.}

\xsection{Chapter XVI.—Watchfulness; The Coming of the Lord}

1. Watch for your life’s sake.\footnote{\par  Or, “over your life;” the clause occurs verbatim in Apostolic Constitutions.} Let not your lamps be quenched, nor your loins unloosed;\footnote{\par  Comp. Luke xii. 35, which is exactly cited in Apostolic Constitutions.} but be ye ready, for ye know not the hour in which our Lord cometh.\footnote{\par  Matt. xxiv. 42.} 2. But often shall ye come together, seeking the things which are befitting to your souls: for the whole time of your faith will not profit you,\footnote{\par  Here Barnabas, iv., furnishes a parallel.} if ye be not made perfect in the last time. 3. For in the last days\footnote{\par  This reference to the last days as present or impending is an evidence of early date; comp. Barnabas, iv., and many passages in the New Testament. The mistake has been in measuring God’s prophetic chronology by our mathematical standard of years.} false prophets and corrupters shall be multiplied, and the sheep shall be turned into wolves, and love shall be turned into hate;\footnote{\par  Comp. Matt. xxiv. 11, 12.} 4. for when lawlessness increaseth, they shall hate and persecute and betray one another,\footnote{\par  Comp. Matt. xxiv. 10.} and then shall appear the world-deceiver\footnote{\par  \textgreek{ὁ κοσμοπλάνος}, found only here and in Apostolic Constitutions, vii. 32. Comp. 2 Thess. ii. 3, 4, 8; Rev. xii. 9.} as the Son of God,\footnote{\par  Not found in Apostolic Constitutions. The expression plainly implies the belief that Jesus Christ was Son of God.} and shall do signs and wonders,\footnote{\par  Comp. Matt. xxiv. 24. The rest of the verse has no parallel.} and the earth shall be delivered into his hands, and he shall do iniquitous things which have never yet come to pass since the beginning. 5. Then shall the creation of men come into the fire of trial,\footnote{\par  Comp. 1 Pet. iv. 12. where \textgreek{πύρωσις} also occurs.} and many shall be made to stumble and shall perish; but they that endure in their faith shall be saved\footnote{\par  Comp. Matt x. 22 and similar passages; none of them directly cited here.} from under the curse itself.\footnote{\par  \textgreek{ὑπ’ αὐτοῦ τοῦ καταθέματος}, “from under the curse itself:” namely, that which has just been described. Bryennios and others render “by the curse Himself;” that is, Christ, whom they were tempted to revile. All other interpretations either rest on textual emendations or are open to grammatical objections. Of the two given above, that of Hall and Napier seems preferable.} 6. And then shall appear the signs of the truth;\footnote{\par  “Truth” might refer to Christ Himself, but the personal advent is spoken of in verse 8; it is better, then, to refer it to the truth respecting the parousia held by the early Christians. For this belief they were mocked, and hence dwelt upon it and the prophecies respecting it. The verse is probably based upon Matt. xxiv. 30, 31; but some find here, as in verse 4, an allusion to Paul’s eschatological statements in the Epistles to the Thessalonians.} first, the sign of an out-spreading\footnote{\par  Professor Hall now prefers to render \textgreek{ἐκπετάσεως}, “outspreading,” instead of “unrolling,” as in his version originally. Hitchcock and Brown, Schaff, and others, prefer “opening;” that is, the apparent o pening in heaven through which the Lord will descend. “Outspreading” is usually explained (so Professor Hall) as meaning the expanded sign of the cross in the heavens, the patristic interpretation of Matt. xxiv. 30. Bryennios and Farrar refer it to the flying forth of the saints to meet the Lord. There are other interpretations based on textual emendations. As the word is very rare, it is difficult to determine the exact sense. “Opening” seems lexically allowable and otherwise free from objection.} in heaven; then the sign of the sound of the trumpet; and the third, the resurrection of the dead; 7. yet not of all, but as it is said: The Lord shall come and all His saints with Him.\footnote{\par  Zech. xiv. 5. This citation is given substantially in Apostolic Constitutions. As here used, it seems to point to the first resurrection. Comp. 1 Thess. iv. 17; 1 Cor. xv. 23; Rev. xx. 5. Probably it is based upon the Pauline eschatology rather than upon that of the Apocalypse. At all events, there is no allusion to the millennial statement of the latter. Since there was in the early Church, in connection with the expectation of the speedy coming of Christ, a marked tendency to Chiliasm, the silence respecting the millennium may indicate that the writer was not acquainted with the Apocalypse. This inference is allowable, however, only on the assumption of the early date of the Teaching.} 8. Then shall the world see the Lord coming upon the clouds of heaven.\footnote{\par  Comp. Matt. xxiv. 30. The conclusion is abrupt, and in Apostolic Constitutions the New-Testament doctrine of future punishment and reward is added. The absence of all reference to the destruction of Jerusalem would indicate that some time had elapsed since that event. An interval of from thirty to sixty years may well be claimed.}

\xsection{Elucidations}

(Thus baptize ye, p. 379.)

If we compare this chapter with the corresponding one in the Apostolic Constitutions, the Teaching seems to me to be a somewhat abridged form of a common original. This being designed for the catechumens, there is an omission of what they are afterwards to know. A form originally drawn up for clergy and people has been very inartificially expurgated for the instruction of young disciples. This appears from the ninth chapter (p. 380), where only certain receptive or responsive forms are given. The liturgy of the Apostolic Constitutions, book viii., embodies what was studiously kept from all but the \textgreek{τέλειοι}, i.e., those “of full age.”

(Concerning Apostles, p. 380, note 16.)

The reference to “apostles,” probably itinerant, in Rev. ii. 2, corresponds with this. There were officers known in the Apostolic day (compare 2 Cor. viii. 23, Greek) as \textgreek{ἀπόστολοι} \textgreek{ἐκκλησιῶν}, for the pseud-apostles of the Apocalypse could not have pretended what they did had it been otherwise. Neither would it have been needful to “try those who said they were apostles,” in that case: the mere assertion of such a pretence would have sufficiently convicted them.

The very childish directions (suited to mere catechumens) given in the text illustrates Rev ii. 2, and is, so far, evidence of the very early origin of the Teaching.

The name apostles was made technical by Christ Himself: “He named them Apostles” (Luke vi. 13). And the word is never used in the loose way which Bishop Lightfoot hazardously suggests, as I must venture to believe.

(Incipient fanaticism, p. 381, note 25.)

Unquestionably, for even in St. Paul’s day his admonitions imply nothing less. See 1 Cor. cap. xiv., passim. But, as in the Introductory Notice\footnote{\par  P. 371, supra.} I hinted my suspicions of incipient Montanism in the Teaching, so I am strengthened in this idea by the learned critic to whose note I venture to append this remark for the purpose of asking a reference to my annotations of Hermas in vol. ii. of this series. May I also ask a reference to the same volume, pp. 4, 5, and 6? The “meal” (note 23, p. 380) of the Teaching is doubtless the Agape, which had been abused at so early a day, that St. Peter\footnote{\par 2 Pet. ii. 13. Compare 1 John iv. 1.} himself was forced to denounce the “false prophets” who polluted this feast of charity.
